% Code for creating comments
\newboolean{showcomments}
\setboolean{showcomments}{true}
\ifthenelse{\boolean{showcomments}}
{\newcommand{\mynote}[2]{
  \fbox{
  \bfseries\sffamily\scriptsize#1}
  {\small
  $\blacktriangleright$
  \textsf{\textcolor{red}{{\em #2}\bf}}
  $\blacktriangleleft$}}
}

% Create custom user defined command for author notes
\newcommand{
  \benoit}[1]{\mynote{Benoit}{#1}
}
\newcommand{
  \martin}[1]{\mynote{Martin}{#1}
}
\newcommand{
  \javier}[1]{\mynote{Javier}{#1}
}
\newcommand{
  \cesar}[1]{\mynote{Cesar}{#1}
}
\newcommand{
  \gluckzhang}[1]{\mynote{Long}{#1}
}

% Create custom user defined command for `TODO` notes
\newcommand*\badge[1]{
  \colorbox{red}{\color{white}#1}
}
\newcommand{\todo}[1]{
  \noindent\textbf{\badge{TODO}} {\color{red}#1}
  \GenericWarning{}{LaTeX Warning: TODO: #1}
}

\newcommand{\review}[1]{\textcolor{blue}{#1}}

% Shortcuts
\newcommand{\ie}{i.e.\@\xspace}
\newcommand{\aka}{a.k.a.\@\xspace}
\newcommand{\eg}{e.g.\@\xspace}
\newcommand{\etal}{et al.\@\xspace}
\newcommand{\wrt}{w.r.t.\@\xspace}
\newcommand{\dapps}{\DH{Apps}\@\xspace}

\newcommand{\chaoseth}{\textsc{ChaosETH}\@\xspace}
\newcommand{\geth}{\textsc{Geth}\@\xspace}
\newcommand{\erigon}{\textsc{Erigon}\@\xspace}
\newcommand{\nethermind}{\textsc{Nethermind}\@\xspace}
\newcommand{\besu}{\textsc{Besu}\@\xspace}
\newcommand{\neth}{\textsc{N-ETH}\@\xspace}

\newcommand*\circled[1]{\tikz[baseline=(char.base)]{
    \node[shape=circle, draw, inner sep=1pt, 
        minimum height=12pt] (char) {#1};}}

 \newcommand{\linebreakand}{%
      \end{@IEEEauthorhalign}
      \hfill\mbox{}\par
      \mbox{}\hfill\begin{@IEEEauthorhalign}
    }

\newcommand{\rqone}{What are the behavioral consequences of unstable execution environments for blockchain nodes?}

\newcommand{\rqtwo}{To what extent do different blockchain node versions exhibit different availability rates under unstable execution environments?}

\newcommand{\rqthree}{To what extent does an N-version blockchain node increase availability compared to a single-version node?}

\usetikzlibrary{shapes.geometric}

\newcommand{\score}{
  \protect\begin{tikzpicture}
      \protect\node[star,
        star points=5,
        star point ratio=2.25,
        draw,
        inner sep=1pt,
        fill=yellow]{};
  \end{tikzpicture}
}

\newcommand{\triup}{
\begin{tikzpicture}
  \node[isosceles triangle,
    isosceles triangle apex angle=50,
    draw,
    rotate=90,
    fill=green,
    inner sep=1pt,minimum size=1pt] (T1)at (0,0){};
  \end{tikzpicture}
  }
\newcommand{\tridown}[1]{
\begin{tikzpicture}
  \node[isosceles triangle,
    isosceles triangle apex angle=50,
    draw,
    rotate=270,
    fill=#1,
    inner sep=1pt,minimum size=1pt] (T1)at (0,0){};
  \end{tikzpicture}
  }
